\documentclass{article}
\usepackage[portuguese]{babel}
\usepackage[letterpaper,top=2cm,bottom=2cm,left=3cm,right=3cm,marginparwidth=1.75cm]{geometry}

\usepackage{amsmath}
\usepackage{graphicx}
\usepackage{amssymb}
\usepackage[colorlinks=true, allcolors=black]{hyperref}

\title{Coisas Úteis que eu aprendi sobre Matemática}
\date{}

\begin{document}
\maketitle

\tableofcontents

\section{O que tem por Aqui?}
Aqui, eu vou anotar praticamente tudo sobre matemática que eu aprendi na vida, que não seja de nível superior ou muito específico. Desde operações básicas, até pré cálculo
\\[5pt]
\textbf{Uma coisa importante:} Às vezes eu explico uma coisa de forma desnecessariamente didática e/ou longa, ou eu explico a mesma coisa de várias formas diferentes. Eu faço isso porque, mesmo que eu esteja explicando da mesma forma, às vezes explicando de um tal jeito ou de outro, faz diferença na compreensão, às vezes um jeito entra mais fácil na cabeça. Ou às vezes sendo desnecessariamente didático a coisa realmente fica mais fácil. O ponto é: Se você entender de primeira um conceito, logo na primeira explicação, nem precisa perder tempo lendo as outras formas de explicar. Elas servem só pra complementar caso esteja difícil entender a primeira explicação 

\section{Como Aprender Matemática?}
Parece que é redundante eu explicar isso né? Mas na realidade, eu acho que é uma das coisas mais importantes, porque antes de eu \textbf{"aprender como aprender"}, eu não conseguia entender nada na matemática. Aqui vão várias dicas que eu aprendi

\subsection{Na Matemática, não é pra generalizar}
Sim, soa estranho. Deixa eu explicar o que eu quero dizer com isso. Sabe quando você aprende uma coisa na matemática, e aquela coisa vale só pra um ou poucos casos, e você percebe que da pra generalizar e/ou percebe que aquele caso que foi ensinado sozinho não cobre a maioria dos casos de problemas que apareceriam? Então, quando acontecer isso, SIGA SÓ O QUE FOI EXPLICADO. Fim. No futuro, isso vai servir pra complementar ou te fazer entender algum outro raciocínio
\\[5pt]
Vou dar um exemplo pra não ficar abstrato: Produtos notáveis e fatoração
\\[5pt]
Você aprende só alguns casos de produtos notáveis, $(a+b)^2$, $a^2-b^2$, trinômio quadrado perfeito etc. Mas aí, bate aquela sensação: "Mas não tá faltando coisa aqui? E se eu quisesse fatorar um polinômio arbitrário?". Isso não vai te levar pra lugar nenhum. Claro, pensar isso de curiosidade só, beleza, agora, achar que você não entendeu como uma coisa funciona só porque ela não é geral é ERRADO. A ideia não é generalizar, e tá tudo certo assim! Veja, quando falamos dos produtos notáveis e fatoração, a maioria das funções e expressões usadas na economia, ciência de dados, cálculo e quase todas as áreas são possíveis de fatorar usando esses casos de produtos notáveis que a gente aprende e/ou combinações deles! A ideia nunca foi você entender como fatora uma coisa maluca tipo $(a - \frac{b^7}{27a+a^2}+ab^{2300}+4a^2) +82ab$, isso é de pouca importância. Isso o computador fazer pra você. O que importa é entender a lógica da fatoração, e os casos mais úteis envolvendo esse tópico
\\[5pt]
Ou seja, era o que eu tinha falado, a ideia não é generalizar, é só aprende o que importa, o que a gente vai usar, e aí, talvez, no futuro, se for útil, generalizar

\subsection{Exercícios não são só exercícios, ELES SÃO PARTE DA TEORIA}
Perceba que eu não estou falando do clássico e batido "Tem que fazer exercícios pra aprender direito, hein", é uma coisa muito mais profunda
\\[5pt]
Veja, se você só aprender um conceito ou uma teoria e não souber aplicar ela direito, você não aprendeu nada! Porque a teoria assume que você vá saber usar o que foi ensinado nos MAIS DIVERSOS CASOS
\\[5pt]
Então, mesmo que você saiba a teoria totalmente, se você nunca aplicou ela, significa que você não aprendeu totalmente.
\\[5pt]
É tipo aprender a usar uma chave de fenda na teoria. Você aprender "Pegue a chave de fenda, encaixe no parafuso e gire para parafusar". Mas percebe que isso não vale de nada? Porque a explicação teórica nunca vai abranger os infinitos tipos de parafuso que tem por aí. Se você só aprender através da teoria, você só saberia parafusar um ou poucos tipos de parafusos ideais
\\[5pt]
Outro exemplo é dirigir. Você vai ensinar a pessoa na teoria, e na prática, claro, trocar marcha, fazer curvas, dar seta, etc. Mas como você pretende ensinar a pessoa de forma teórica a dirigir em infinitos tipos diferentes de condições e ambientes? É impossível. Percebe como a prática é parte da teoria? A pessoa só consegue aplicar a teoria se ela aprender na prática. É esquisito, mas faz sentido, "A teoria só consegue ser ensinada na prática". Vale pra matemática e pra muitas coisas na vida

\section{Assuntos Essenciais de Matemática Básica}

\subsection{Coisas pra Refrescar a Memória}
Eu vou deixar anotado aqui umas coisas muito simples que eu não preciso explicar a fundo, e nem explicar pra falar a verdade, são muito fundamentais. Tô anotando aqui só pra refrescar a memória 

\subsubsection{Multiplicação com números Fracionários}
Quando aparece multiplicação assim, é só contar o número de casas decimais depois da vírgula em TODOS os números que foram multiplicados. Aí multiplica normal. Depois, você vai contar de trás pra frente no resultado da multiplicação o mesmo tanto de casas decimais que você tinha contado no começo. Pronto, é isso

\subsubsection{Divisão com Dividendo Menor que o Divisor}
Quando isso acontece, adiciona um zero no dividendo e um 0, no quociente. Caso mesmo depois de fazer isso o dividendo ainda seja menor que o divisor, adiciona mais um zero no quociente. Caso ainda seja menor, vai adicionando zero. Pronto, quando o dividendo for maior que o divisor é só dividir normal

\subsubsection{Divisão com número fracionário}
Coloque zeros até que os números deixem de ser fracionários. Seja o divisor, o dividendo ou ambos. Quando não tiver mais vírgula, faça a divisão normalmente.

\subsection{Divisão}
\textbf{A TÃO TEMIDA DIVISÃO.} Ah... A operação de divisão... Pensar tudo que ela representou pra sua vida é assustador. Os problemas de matemática do 5° ano que te destruíram, a soma de frações que te atormentou até os 20 anos, as razões e proporções que você demorou uma vida pra entender... Lembrar da divisão é lembrar da sua defasagem escolar, do ensino fundamental 1 mal feito, do fundamental 2 quase todo ausente e do ensino médio, que foi na pandemia, então você não aprendeu nada. Mas não se preocupa, tenta esquecer esse monte de merda. Eu te ajudo com a divisão
\begin{itemize}
    \item \textbf{O QUE É A DIVISÃO?}\\[5pt] Isso parece meio óbvio, mas é bom definir formalmente
    \\[5pt]
    A divisão é uma operação divide um valor em tantas partes iguais. O resultado da divisão é o tamanho de cada uma dessas partes iguais. Tipo, 10÷2 = 5. O resultado é 5 porque o tamanho de cada parte que foi dividida é 5
    \\[5pt]
    De forma mais fundamental e complicada, tipo, de forma elementar, a divisão representa uma sequência de subtrações. É fácil de entender quando você pensa na multiplicação, do mesmo jeito que a multiplicação representa uma sequência de somas, a divisão, sendo o oposto, representa uma sequência de subtrações
    \\[5pt]
    \textbf{IMPORTANTE:} Em quase todos os casos de divisão do dia a dia, e na maioria dos casos envolvendo problemas com divisão, não tem vantagem nenhuma pensar na divisão como sendo uma sequência de subtrações, isso só vai complicar. \textbf{E quando isso é útil?} Na hora de estudar construção dos números, lógica, na hora de provar teoremas etc

    \item \textbf{O QUE A DIVISÃO REPRESENTA}\\[5pt] Agora vem a parte meio complicada. A divisão representa mais do que um tipo de operação. Calma, vamos com calma
    \\[5pt]
    A parte óbvia primeiro:
    \\[5pt]
    \textbf{A divisão representa uma repartição.} Fácil, isso é o que qualquer um que fez ensino fundamental sabe. Eu nem precisaria explicar porque realmente qualquer um sabe. Mas vamos lá, quando você tem um total A e quer repartir pra um outro total B de C partes iguais, você usa a divisão. Tipo, tenho 27 balas e quero dividir pra 9 pessoas, com quantas balas cada pessoa vai ficar? Simplesmente divide 27 por 9 e você vai obter a quantidade de balas que cada pessoa vai ter depois da repartição. Isso da 3, então cada pessoa vai ficar com 3 balas. Totalmente coisa do ensino fundamental, mas ainda assim, é um dos usos mais comuns da divisão 
    \\[5pt]
    Agora, as partes não óbvias. Sim. No plural, porque tem bem mais de uma, afinal a divisão representa várias coisas.
    \\[5pt]
    \textbf{Quantas vezes uma coisa cabe dentro da outra?} A divisão pode representar quantas vezes um valor "cabe" dentro de outro valor. É meio confuso pensar assim, mas quando você pensa que a divisão é o oposto da multiplicação fica mais fácil entender
    \\[5pt]
    \textbf{Exemplo:} Eu tenho uma caixa do tipo $A$ com $1000m^3$ de volume, e caixas menores, caixas do tipo $B$, cada uma com $125m^3$ de volume. Quantas caixas do tipo $B$ cabem dentro de uma caixa do tipo $A$? E quantos porcento do volume de uma caixa do tipo $A$ cabem dentro de uma caixa do tipo $B$?
    \\[5pt]
    Pra resolver esse problema, você usa divisão. Mas por que você usa a divisão? Porque você aplica a lógica que eu falei, \textbf{``Quantas vezes uma quantidade cabe dentro da outra?"}, e sabemos que quando usamos essa lógica, trata-se de uma divisão. Você vai ter que descobrir quantas partes de $125m^3$ cabem dentro de $1000m^3$. Pra fazer isso, é só dividir 1000 por 125, que da 8. Sabe por que? Porque usamos a divisão quando queremos descobrir quantas vezes uma quantidade se repete para formar outra — e a divisão é a operação que desfaz a multiplicação usada pra repetir essa quantidade
    \\[5pt]
    \textbf{ENTÃO, LEMBRE-SE:} Você usa divisão porque ela responde exatamente à pergunta “quantas vezes uma quantidade cabe dentro da outra?”
    \\[5pt]
    \textbf{Quanto uma quantidade representa da outra?} Essa é outra utilidade da divisão! Por exemplo, você quer saber quanto 2 representa de 10. Pra saber quantas vezes 2 cabe dentro do 10 a gente já saberia fazer, só dividir 10 por 2, agora, e pra saber quanto o 2 representa do 10? Fácil! Só dividir o 2 por 10, o inverso do que a gente faria pra saber quantas vezes o 2 cabe dentro do 10

\end{itemize}

\subsection{Problemas Envolvendo Divisão (IMPORTANTE!!!)}
Nessa parte, eu vou deixar uns problemas EXTREMAMENTE úteis envolvendo razão, proporção, divisão, e os caralho tudo

\begin{itemize}
    \item Um menino tem 10,2 anos de idade. Quantos anos, meses e dias esse menino tem?
    \\[5pt]
    Simples, 10 anos eles já tem completo, então só resta descobrir quantos anos 0,2 anos representa. 1 ano tem 12 meses, então, só multiplicar 0,2 por 12 pra saber quantos anos 0,2 anos representa! Às vezes fazer isso pode não parecer muito natural, porque quando o valor é menor do que 1, fica estranho. Mas pensa, se eu falo pra descobrir quantos meses são 2 anos, sabendo que 1 ano = 12 meses, é só multiplicar 12 por 2. No caso, para descobrir quantos anos são 0,2 anos, mesma coisa! Você simplesmente multiplica. Temos que $0,2 \times 12 = 2,4$, então o menino tem 10 anos e 2,4 meses. Pra descobrir quantos anos, meses e dias ele tem, fazemos o mesmo esquema de multiplicar. 2 meses já foram completos, então vamos descobrir quantos dias são 0,4 meses. 1 mês = 30 dias então 0,4 meses = $0,4 \times30 = 12$. Pronto! O menino tem 10 anos 2 meses e 12 dias de vida! O problema é bem óbvio e fácil de resolver, eu nem precisaria explicar de forma tão detalhada. O que complica mesmo é a quantidade de meses (0,2) ser um valor menor que 1, isso deixa um pouco contraintuitivo de pensar sobre, mas com calma da certo!

    \item 6 pés são 183cm. Quantos metros tem em 287 pés?
    \\[5pt]
    Absurdamente fácil, divide 183 por 6 pra saber quantos cm vale 1 pé, afinal, a divisão vai distribuir igualmente o 183 por 6, o resultado é 30,5. Então 1 pé vale 30,5cm. Se sabemos disso, então a gente multiplica 287 por 30,5 pra saber quantos pra saber quantos cm essa quantidade de pés tem. O resultado é 8.753,5cm. Agora, precisamos passsar isso pra metros. Como a gente sabe que 1 metro = 100cm, a gente divide 8.753,5 por 100 pra saber quantos metros cabem dentro de 8.753,5. O resultado são 87,535m. Pronto!

    \item PROBLEMA DA BARRA DE LOADING: Uma barra de loading A carregou 37\%, de 100\%. Temos outra barra de loading B, arbitrária, que inicia em 68 e termina 487. Quanto o 37\% da barra de loading $A$ equivale ao da $B$?
    \\[5pt]
    Aqui, temos um caso que da pra usar a lógica do ``raciocínio por unidade", afinal, essa lógica sempre é aplicável, mas é bem mais fácil e até mais logicamente intuitivo, simplesmente multiplicar direto a porcentagem de uma barra pelo total da outra. É natural, você sabe que se você multiplicar a \% da barra $A$ pelo total da outra, o resultado vai ser a \% equivalente da barra $B$. Só toma cuidado, porque a barra $B$ começa em 68, então o total da barra $B$ é 487 - 68, que é 419. Agora sim, temos que o tamanho total da barra $B$ é 419, então só multiplicar por 37\%, e depois adicionar 68 de novo, afinal, ela começa 68 unidades a frente do zero. $37\% \times 419 = 155,03 \Rightarrow 155,03 + 68 \approx 223$. Pronto! Temos que 37\% de 100\% da barra de loading $A$ equivale a 223\% de 487\% da barra de loading $B$
    \\[5pt]
    \textbf{Agora, a forma não recomendada}. Vamos por lógica agora, mesmo não sendo recomendado por ser, justamente, meio desncessário e até confuso ir por lógica nesse caso. Primeiro, você tem que descobrir quantas carregadas tem pra cada unidade do total da barra $A$, e depois multiplicar essa quantidade pelo total da barra $B$. Ou seja, primeiro faz $37\div100$, que resulta em 0,37. Depois, multiplica 0,37 por 419, que resulta em 155,03, igualzinho a resolução anterior . Agora é só somar 68, pronto, 223. E por que eu falei que é confuso? Porque a porcentagem confunde. Como as unidades carregadas (37\%) e o total (100\%) são porcentagens, pode confundir, porque se você fizer $37\%\div100$, vai dar errado. Ou você usa ambos como porcentagem, fazendo $37\%\div100\%$, ou você faz do jeito que eu fiz, ignorando o fato desses valores serem porcentagens, e simplesmente tratando eles como unidades, afinal, nesse caso em que estamos indo por lógica, não importa se o 37 e o 100 são porcertagens, o que importa é saber quantas unidades de uma quantidade tem pra cada unidade do total da outra
    \\[5pt]
    FORMA AINDA MAIS ALTERNATIVA: Ainda tem outros jeitos de resolver que seriam considerados ``normais". Olha esse jeito. A barra $B$ vai de $68$ até $487$, então o tamanho total dela é: $487 - 68 = 419$. Se $100\%$ da barra $B = 419$ unidades, então $1\%$ da barra $B = \frac{419}{100} = 4{,}19$ unidades. Esse é o raciocínio por unidade: descobrir quanto vale $1\%$. Agora, se $1\%$ vale $4{,}19$, então $37\%$ vale $37 \times 4{,}19 = 155{,}03$. Como a barra começa em $68$, o ponto equivalente será: $68 + 155{,}03 = 223{,}03$. Meio diferentão, mas funciona, e você lida diretamente com porcentagens, diferente da resolução anterior, que lida puramente com a relação entre unidades de forma arbitrária

    \item Converta 7800 dias para anos \\[5pt] Muito simples, certo? Só dividir por 7 pra converter pra semanas, depois que estiver em semanas, só dividir por 4 pra passar pra meses e depois dividir por 12 pra passar pra anos. Certo? Errado! Mas por que errado?! Errado porque 1 mês não tem exatamente 4 semanas, são 4,5 semanas aproximadamente. Agora sim, se você dividir por 4,5 quando for passar pra meses, da \textit{mais certo}, mas ainda da meio errado. Perceba uma coisa, sempre vai ter um erro de aproximação pra passar de semanas pra meses, porque 1 mês não é exatamente 4,5 semanas. Isso acontece em muitos, muitos casos com conversões. Acontece na maioria dos casos pra falar a verdade. Então esse problema virá a nos ensinar uma coisa essencial: TENTE REALIZAR O MÍNIMO DE CONVERSÕEs POSSÍVEIS. Sabemos que 1 mês são 30 dias, certo? Então por que passar pra semanas e depois pra meses? Já converte direto de dias pra meses! SEMPRE faça isso. Evite ao máximo conversõe, faça apenas as necessárias. Então, pro nosso problema, vamos dividir 7800 por 30, e depois por 12, pra passarmos pra meses e pra anos respectivamente. Então, temos que 7800 dias são aproximadamente 21,6 anos. Certo? ERRADO!!!!!!!!!!!!!!!!!!!!!!!!! Sabemos que 1 ano são 365 dias então POR QUE passar pra meses??????? Converte direto pra anos. $7800\div365$, temos que 7800 dias são aproximadamente 21,4 anos. Essa é a melhor aproximação que conseguimos fazer
    
\end{itemize}

\subsection{Frações}
Agora, chegou a vez das tão temidas frações. O que elas são? Qual a diferença de uma fração e uma divisão comum? Vamos descobrir!

\begin{itemize}
    \item \textbf{O Conceito de Fração}\\[5pt] Podemos pensar em uma fração como uma forma de representar uma \textit{``divisão regulável"} ou uma \textit{divisão mais complexa}. Só que é interessante entender o que ela realmente mostra/representa. Quando você escreve uma fração, tipo $\frac{3}{4}$, o número de baixo (o denominador) diz em quantas partes iguais um total foi dividido, e o número de cima (o numerador) diz quantas dessas partes você ``pegou". Como fica meio feio falar ``pegou" em um contexto matemático, falamos ``tomou". Então $\frac{3}{4}$ significa “você pegou/tomou três partes de um total dividido em quatro partes iguais”, ou seja, você pegou um total, dividiu em 4 partes, e tomou 3 dessas partes
    \\[5pt]
    O mais importante é perceber que uma fração é só um jeito de escrever uma divisão que ainda não foi feita. Tipo, $\frac{3}{4}$ é o mesmo que $3 \div 4$, mas em forma “guardada”. É por isso que dá pra operar frações sem precisar transformar em decimal — elas mantêm a divisão “em espera”.

    \item \textbf{DIVISÃO VS FRAÇÃO}\\[5pt] Apesar de ambos os conceitos de divisão e fração envolverem divisões, se compararmos uma divisão com uma fração, a principal diferença é: \textit{A divisão é uma operação, enquanto a fração é um conceito que envolve divisão}
    \\[5pt]
    O que entendemos disso: A divisão simplesmente divide um total em tantas partes iguais. A fração, por outro lado, é um conceito mais complexo, ela indica em quantas partes iguais um total foi dividido, e quantas dessas partes iguais foram tomadas/escolhidas/consideradas
\end{itemize}

\subsection{Múltiplos e Divisores}

\subsubsection{Múltiplo}
Um múltiplo é um número $a$ que aparece na taboada de $b$. Ou seja,

\subsubsection{Divisor}

\subsection{A Diferença Entre Múltiplo e Divisor}
Eles podem parecer conceitos idênticos, mas perceba a diferença

\subsection{A parte pelo todo, o todo pela parte}
Isso é a lógica por trás da razão e da proporção!

\subsection{Razão}
Vamos falar sobre a razão, e também brevemente sobre a famosa proporção, que na maioria das vezes é chamada de regra de três

\begin{itemize}
    \item \textbf{O que é uma razão?}\\[5pt] Uma razão é uma fração (ou uma divisão) que mostra quantas vezes uma quantidade cabe dentro da outra, relacionando assim, 2 quantidades
    \\[5pt]
    Sim, eu sei uma operação de divisão comum já mostra quantas vezes uma quantidade cabe dentro de outra. A diferença é que a razão não é uma operação, nem um tipo de conta que você faz, ela é uma constante, ou uma relação, ou seja, \textbf{a razão é o resultado de uma divisão entre duas quantidades que você quer comparar}
    \\[5pt]
    Você também pode pensar assim: Uma razão gera um valor que te diz: "a cada tantas coisas x, você terá tantas coisas y". Perceba que isso é o mesmo que dizer "quantas vezes x cabe em y”, já que ambas as lógicas comparam as duas quantidades pela mesma relação de divisão. Porém, às vezes fica mais fácil pensar desse jeito. Calma, já exemplifico
    \\[5pt]
    Vamos pensar em um exemplo prático: Se falarem "1 dólar = 2 reais", é como se 1 dólar fosse um pacote com 2 reais dentro. Então, se tivermos 5 dólares, e quisermos saber quantos reais isso vale, é só pensar que é como se tivéssemos 5 pacotes com 2 reais dentro, logo 10 reais. Agora, se quisessemos saber quantos dólares vale 10 reais, você teria que pensar quantos pacotes de 2 reais cabem dentro desse valor 10, que seriam 5 pacotes, logo 5 dólares. Agora que vem a sacada! Perceba que dizer ``A cada tantos dólares da pra formar tantos pacotes de reais" equivale a dizer ``A cada tantos dólares teremos tantos reais"
    \\[5pt]
    Então, concluímos que dizer ``a cada tantas unidades de X teremos tantas unidades de Y" é equivalente a falar: \textit{A cada tantas unidades de X, da pra colocar quantos "pacotes" de Y dentro?}

    \item \textbf{Exemplo de Razão:} $\frac{4}{5}$. Isso representa uma razão pois ela mostra quantas vezes uma quantidade cabe dentro de outra. Vamos dizer que a quantidade 4 é X e a quantidade 5 é Y. Perceba! Se dividirmos 4 por 5 obtemos quantas unidade de X temos para cada unidade de Y, que seria 0,8. Isso já indica a razão. Cada 0,8 unidades de X temos 1 unidades de Y. Ou seja, 0,8 é nossa razão de proporção. Ou você poderia pensar assim: Essa divisão 4/5 está indicando quantas vezes 5 cabe dentro de 4. 5 cabe dentro de 4 0,8 vezes, então para cada 1 do denominador (5), você tem 0,8 do numerador (4). Logo, temos uma porporção de 1:0,8 ou 4:5 (vamos ver essa notação dos dois pontos mais pra frente)
    
    \item \textbf{Porque a razão não é uma divisão/fração comum}\\[5pt] Isso mesmo que você ouviu. Eu faço questão de reafirmar. A razão não é uma divisão/fração comum. Sim, claro, ela é uma divisão/fração, mas ela tem uma utilidade bem específica na matemática. A divisão é simplesmente uma operação, e a razão é a forma de usar essa operação para expressar a relação entre duas quantidades, gerando uma constante
    \\[5pt]
    \textbf{``Toda razão é uma divisão, mas nem toda divisão representa uma razão"}
    \\[5pt]
    A razão é um uso específico da divisão pra relacionar 2 grandezas. Diferente de uma divisão arbitrária, que pode ter a intenção de simplesmente separar partes iguais, ou qualquer outra intenção arbitrária. Fazendo uma analogia:
    \\[5pt]
    A divisão é como uma palavra - algo simples e isolado.
    \\
    A razão é como uma frase, que usa palavras pra transmitir uma ideia completa.
    \\
    A razão nasce da divisão, mas expressa uma relação, não só um resultado.

    \item \textbf{A Razão na Forma de Notação}\\[5pt] Sabe o famoso 1:1 ou 1:3 ou 4:10? Isso são razões. Essa é a notação de razão. Se pegarmos uma razão de 1:2, ela te diz que para cada 1 unidade da primeira coisa, você tem 2 unidades da segunda. É literalmente isso: "A cada 1, correspondem 2.”
    \\[5pt]
    Então, é uma razão escrita no formato de proporção. Ela diz como duas quantidades se comparam entre si. Se você quiser transformar 1:2 em número, é só dividir: 1/2 = 0,5. Então 1:2 também significa “uma coisa é metade da outra”. Ou seja, é a mesma lógica de "uma coisa cabe dentro da outra tantas vezes"
    \\[5pt]
    Quando pensamos na relação de razão, normalmente comparamos 2 quantidades usando uma unidade de seja lá qual quantidade como referência. Por isso que a maiorias razões são 1:X, "1 pra X", por exemplo escala de mapa, é sempre 1:4 ou 1:1000 ou alguma relação/escala usando 1 como referência. Fazemos isso porque é intuitivo, se sabemos quantas unidades de tal quantidade temos por unidade singular da outra quantidade, fica fácil aplicar pra outras quantidades
    \\[5pt]
    Exemplo: X e Y estão em razão de 1:4. Se para cada 1 unidade de uma quantidade (X) temos 4 unidades da outra (Y), fica fácil entender que, para obter a quantidade de X em relação a Y, é só multiplicar, ou dividir. Se você tem 20 unidades de X, só multiplicar por 4 e você teria a quantidade de Y, 80 unidades. Se você tem 200 unidades de Y, fica fácil saber que é só dividir por 4 e você obterá a quantidade de X, que seria 50
    
\end{itemize}
    
\subsection{Proporção}
De forma bem direta: A proporção é uma equivalência entre 2 ou mais razões
\\[5pt]
De forma intuitiva, a proporção funciona assim: 
\\[5pt]
O QUE COMPLICA NA HORA DE CALCULAR A PROPORÇÃO É QUE... Existem muitas formas de resolver proporções. Muitas lógicas, raciocínios e métodos. \textbf{Pra ficar bom em proporção, o segredo é dominar esses métodos}

\begin{itemize}
    \item \textbf{As diferentes Lógicas por trás do Cálculo da Proporção}\\[5pt] Tem várias formas de entender como funciona a proporção, vamos ver cada uma delas, e ver em qual caso usar cada uma! 
    
    \item \textbf{Lógica do ``Raciocínio por unidade"} \\[5pt] Vamos chamar esse método de ``Proporção por Lógica de Unidade"
    \\[5pt]
    Ele consiste em: Descobre quantas unidades de uma quantidade tem pra cada unidade de outra quantidade, e depois multiplica essa razão por uma nova quantidade
    \\[5pt]
    \textbf{QUANDO USAR?} Esse método é bem lógico e bem útil pra problemas que envolvem mais um uso de lógica, pra proporções que envolvem quantidades físicas, pra proporções em que claramente existem comparações entre unidades ou proporções nem tão óbvias e/ou solúveis por outros métodos
    \\[5pt]
    \textbf{Exemplo:} Um carro percorre 180 km com 12 litros de gasolina. Quantos quilômetros ele percorre com 1 litro? E com 20 litros?
    \\[5pt]
    Primeiro, descobre quanto ele anda com 1 litro, é isso que o raciocínio por unidade faz: $\frac{180}{12}=15$. Então 1 litro $\rightarrow$ 15km. Agora fica fácil, $15\times20=300$. Aqui o raciocínio por unidade é ideal, porque você primeiro acha o valor de 1 unidade (no caso, “por litro”) e depois multiplica pra qualquer quantidade desejada.

    \item \textbf{Lógica da Parte Pelo Todo} \\[5pt] Esse método consiste em pensar do seguinte jeito: ``Vou descobrir quanto um valor representa do outro usando razão, e depois que eu tiver a razão de proporção, vou multiplicar pelo novo valor, assim obtendo um valor proporcional." Esse método já é muito próximo da regra de 3, que consiste em igualar razões. Nesse caso, você tá fazendo uma regra de 3 de forma mais direta, e mais lógica. Você só pensa ``Quanto A representa de B?" E depois multiplica. Mais direto que a Proporção por Raciocínio de Unidade e menos direto que a regra de 3
    \\[5pt]
    \textbf{QUANDO USAR?} Quando você tem proporções em que não é fácil identificar como os valores unitários das quantidades se comportam, e/ou quando as unidades das quantidades não tem uma relação clara
    \\[5pt]
    \textbf{Exemplo:} Um tanque de água cheio comporta 60 litros. Se ele está com 24 litros, e outro tanque é duas vezes maior, com 120 litros quando cheio, quantos litros teria esse segundo tanque se estivesse cheio na mesma proporção que o primeiro?
    \\[5pt]
    Primeiro, usamos o raciocínio parte pelo todo pra descobrir qual parte do total o primeiro tanque está cheio. $\frac{24}{60}=0,4$. Isso significa que o primeiro tanque está 0,4 (ou 40\%) cheio. Passo 2, aplicar essa mesma proporção ao segundo tanque: O segundo tanque é duas vezes maior (120 litros), e vai estar cheio na mesma proporção, então: $0,4\times120=48$. Resposta final: O segundo tanque tem 48 litros de água.
    \\[5pt]
    Aqui o “parte pelo todo” serviu pra descobrir a proporção da quantidade de água pelo tamanho (24/60), e depois multiplicamos pra aplicar a mesma proporção num novo total. É bem o raciocínio “parte ÷ todo × novo todo”.
    
    \item \textbf{Proporção Composta}\\[5pt] Quando temos mais de uma grandeza, ou seja, mais de uma quantidade pra entendermos proporções e relações, temos uma proporção composta
    \item \textbf{Razão de Proporção}\\[5pt] Existem vários casos em que é mais fácil ir no ``modo automático", ou seja, não precisa usar a ``Proporção por raciocínio de unidade" de forma crua e pura, afinal, você já sabe como ele funciona, e sabe que ele é aplicável no caso, e que também não é necessário usar regra de três, por motivos de conveniência. Nesses casos, podemos dividir um valor pelo outro pra descobrir a quantidade de X por unidade de Y, e multiplicar. O resultado dessa razão é a 

\end{itemize}

\subsection{Exercícios Envolvendo Razão e Proporção}

\begin{itemize}
    
    \item \textbf{ANTES DE TUDO, UMA COISA MUITO IMPORTANTE:} Cuidado pra não confundir problemas envolvendo proporção com problemas que são simplesmente "quantas vezes x cabe dentro de y?", esses problemas usam simplesmente razão. Ou seja, esses problemas podem ser resolvidos apenas com uma divisão
    \\[5pt]
    exemplo: você andou 11,3km em 2h03min. quanto tempo demorou pra você percorrer 1km? 
    \\[5pt]
    Nesse caso, você só precisa dividir 123 minutos (2h03min) por 11,3. não precisa usar regra de 3. sabe por que? porque você quer descobrir em quanto tempo você andou em 1km, você quer saber em quanto tempo você andou em 1 unidade de km, aí não precisa da regra de 3 né, porque pensa comigo, só dividir 123 minutos por 11,3 já vai te retornar quantos minutos foi por km
    \\[5pt]
    \textit{Mesma coisa vale pra esse exemplo:}
    \\[5pt]
    1 lira turca = 0,024 dólares. quantas liras turcas valem 1 dólar?
    \\[5pt]
    Presta atenção, como ele só quer saber quantas liras turcas vale 1 unidade de dólar, não precisa de regra de três! eu sei que é óbvio, mas as vezes você vai no "modo automático" e acaba pensando errado. nesse caso é só dividir 1 por 0,024. pra descobrir quantas vezes 0,024 "cabe" dentro do 1, ou seja, quantas liras turcas "cabem" em 1 dólar. aí pronto, descobriu quanto 1 dólar vale em liras turcas
    \\[5pt]
    \textbf{ATENÇÃO: RESOLVA OS PROBLEMAS A SEGUIR SEM USAR REGRA DE 3, PRA VOCÊ REALMENTE APRENDER COMO FUNCIONA A PROPORÇÃO}

    \item Um celular tem espaço para gravar 3 horas de vídeo em qualidade normal ou 2 horas em alta qualidade. Se já foram gravadas 2 horas de vídeo em qualidade normal, qual é o tempo que resta para gravar vídeos em alta qualidade?

    Tem várias formas de resolver esse problema. A resolução oficial da OBMEP é:
    \\[5pt]
    Três horas (180 minutos) em qualidade normal correspondem a 2 horas (120 minutos em qualidade alta); dividindo proporcionalmente por 3, concluímos que 180 ÷ 3 = 60 minutos em qualidade normal, correspondem a 120 ÷ 3 = 40 minutos em alta qualidade.

    \item Em uma prova de 15 questões, eu acertei 12 questões. De 0 a 10, quanto eu tirei nessa prova?
    
    Isso é um problema de proporção bem simples. Quando o enunciado fala “de 0 a 10”, ele quer transformar essa relação parte/todo numa escala nova, que vai de 0 até 10, ou seja, é a mesma proporção, só que com um total diferente (antes o total era 15, agora o total é 10).
    \\[5pt]
    Primeiro: Qual a forma mais natural de resolver esse problema? Qual forma de calcular proporção usar?
    \\[5pt]
    Pelo fato de ser uma proporção bem simples e direta, a gente não precisa usar um raciocínio muito profundo. Ou seja, \textbf{o método de "Proporção pelo Raciocínio de unidade" é desnecessariamente complexo aqui.} O que da pra perceber é que, usando o usando raciocínio da parte pelo todo, conseguiremos saber quanto 12 representa de 15, e aí já temos a razão de proporção. $12\div15=0,8$, aí naturalmente só multiplicar po 10 e temos o resultado, que é uma nota 8 de 10. Perceba uma coisa, isso que nós fizemos, apesar de usar o raciocínio da parte pelo todo, JÁ É A PROPORÇÃO! Quando você faz parte ÷ todo = 12 ÷ 15 = 0,8, você tá usando a razão parte/todo, isso mostra a fração do total que você acertou (ou seja, 80\%). Quando depois você multiplica por 10, pra transformar essa fração numa nota, aí você cria uma proporção: está projetando a mesma razão numa escala diferente (de 0 a 10)
    \\[5pt]
    Então, respondendo qual método é o mais natural de usar: É a proporção simples mesmo por regra de 3 OU simplesmente usando parte pelo todo e depois multiplicand por 10

    \item O Brasil tem uma população de 211 milhões. Existem 323 mil pessoas com o sobrenome ``Cordeiro" no Brasil. Portugal tem uma população de 10 milhões (aproximadamente). Existem 26 mil pessoas com o sobrenome ``Cordeiro" em Portugal. Qual país tem uma maior densidade de pessoas com o sobrenome Cordeiro?
    \\[5pt]
    \textbf{Dica:} Esse problema envolve apenas razão/divisão simples!
    \\[5pt]
    Pra calcular isso é fácil, vamos descobrir qual país tem mais pessoas com o sobrenome Cordeiro pra cada unidade de habitante! Fazemos isso usando a divisão, como já sabemos, afinal ela divide igualmente o valor de pessoas com o sobrenome Cordeiro para cada unidade de habitante, então obtemos a quantidade de pessoas com sobrenome Cordeiro por habitante. No Brasil, o cálculo seria $323.000\div211.000.000$. Podemos simplificar essa expressão dividindo ambos os valores por 1000. Temos $323\div211.000$, que resulta em aprox. 0,00153. Então no Brasil, pra cada habitante, temos 0,00153 pessoas com sobrenome Cordeiro. Agora vamos calcular em Portugal, o cálculo é $26.000\div10.000.000$. Também podemos simplificar, obtendo $26\div10.000$, que é 0,0026. Ou seja, em Portugal, pra cada habitante temos 0,0026 pessoas com o sobrenome Cordeiro. Se no Brasil, pra cada habitante tem 0,0015 pessoas com sobrenome Cordeiro, e em Portugal pra cada habitante tem 0,0026 pessoas com esse sobrenome, então quer dizer que em Portugal tem mais pessoas com sobrenome Cordeiro por habitante do que no Brasil, logo, Portugal tem uma maior densidade de pessoas com sobrenome Cordeiro! 
    
    \item \textbf{Problema dos Livros:} Eu li 7 páginas de um livro de 30. Quantas páginas eu teria que ler de um livro de 200 pra que o tanto que eu li no livro de 200 representasse o mesmo tanto do total que eu li do livro de 30?
    \\[5pt]
    \textbf{RESPOSTA:} Eu teria que dividir o 7 por 30 pra descobrir quantas páginas lidas tem pra cada unidade do total de páginas, ou seja, quantas páginas lidas tem em relação ao total de páginas, e depois multiplicar o resultado dessa divisão por 200, que é o novo total de páginas. É sempre essa lógica do "descobre quanto tem pra 1 unidade de tal total e depois multiplica pela unidade de outro total"

    \item \textbf{Problema de Portugal:} Vamos pegar um voo pra portugal partindo de São Paulo, 18:00h da tarde. o voo dura 9 horas. portugal tem um fuso horário de GMT+1, e São Paulo de GMT-3. ou seja, o horário de Portugal está 4 horas a frente do horário de São Paulo
    \\[5pt]
    Caso quiséssemos calcular que horas seria em Portugal quando a gente chegasse, seria fácil. Pensa, saímos do Brasil 18:00, são 9h de voo, e portugal está 4 horas adiantado em relação a São Paulo. Absurdamente simples, só somar! 18:00h da tarde + 9h de voo, dá 03:00 horas da manhã MAS como o horário lá é 4 horas adiantado, só somar +4! Quando a gente chegasse, seria 07:00h da manhã
    \\[5pt]
    Agora, vou complicar o problema. E se a gente quisesse generalizar e calcular qual seria o horário local de cada ponto da viagem? Tipo, passou 5 horas de voo, a gente quer descobrir que horas é no fuso horário local naquele ponto do mundo (que provavelmente seria no meio do oceano atlântico). Pra fazer isso é bem simples na verdade, pensa comigo: De São Paulo até Portugal da 9 horas de voo, certo? E de São Paulo até Portugal o fuso horário avança 4 horas. Então, a gente precisa descobrir quanto o fuso horário avança por hora do tempo total de viagem. Pra isso é muito simples, divide 4 por 9! Aí, a gente sabe quanto o fuso horário avança a cada hora da viagem, afinal, dividindo a gente distribuiu o 4 pelas 9 horas. Aí, feito isso, é só multiplicar pelo tempo de viagem que a gente escolher, e somar com o tempo da viagem escolhido. No nosso caso queríamos saber pra 5 horas de viagem quanto o fuso avançaria, então seria 4/9 * 5 + 5, isso da aprox 7,2h. Se partimos às 18:00, então depois de 5 horas de viagem, naquele ponto no meio do oceano atlântico, o horário local naquele fuso horário aleatório do meio do atlântico seria 18:00h + 7,2h, ou seja mais ou menos 1 da manhã
    \\[5pt]
    E sim, daria pra fazer uma proporção também se quisesse. Se em 9 horas de voo o fuso horário avança 4 horas, em quantas horas ele avança pra tal X horas de viagem? No nosso caso, a gente queria saber quanto avançou pra 5 horas da viagem, seria: 9 está para 4 assim como 5 está para X
    \\[5pt]
    EXTRA: Modelar o problema matematicamente através de funções!
    \\[5pt]
    Sim, poderíamos criar uma função que tem como domínio o conjunto com os valores do tempo de viagem, e como imagem, teríamos um conjunto com os valores do horário atual naquele ponto do mundo
    \\[5pt]
    Ou seja, a função exibiria pra cada valor de X, representando o tempo de viagem transcorrido desde a partida, um valor de Y, representando o fuso horário naquele ponto do mundo pra aquele horário da viagem
    \\[5pt]
    Pra modelarmos essa função, é fácil! Pense: pra descobrirmos o horário em um dado momento da viagem, precisamos da constante que mostra quanto o fuso horário avança por hora do tempo total da viagem, essa constante é 4/9, lembra? E para descobrirmos quanto o fuso horário avança pro horário escolhido, basta multiplicar essa constante pelo horário novo, afinal, multiplicando cada unidade que mostra quanto tempo o fuso horário avança por hora, por um escolhido horário, temos quanto o fuso horário avançou pro tal horário. Depois, somamos o tempo de viagem escolhido. Precisamos acrescentar o tempo de viagem até o momento que escolhemos porque esse tempo já passou, em relação ao horário antigo, ou a qualquer horário, se passou X horas de viagem, óbvio, temos que somar
    \\[5pt]
    Dito isso, temos que:
    \[
    f(x)=\frac{4}{9} \cdot x+x
    \]
    Dando uma arrumada na fração pra ficar mais compacto: \[f(x)=\frac{4x}{9}+x\]
    O domínio dessa função são valores de 0 a 9 sem incluir o zero nem o nove. O domínio é esse porque o tempo total de viagem até Portugal é 9 horas! Então temos:
    \[
    f:(0,9)\to\mathbb{R},\quad f(x)=\frac{4x}{9}+x
    \]
    Pronto, agora você pode tratar esse problema de forma algébrica :D

    \item \textbf{Qual país cresceu mais em relação ao seu tamanho?}
    \\[5pt]
    Queremos saber qual país cresceu mais percentualmente nos dados intervalos de tempo. O Brasil, que em 2010 tinha 190.700.000 habitantes e em 2020 tinha 211.800.000, ou a espanha, que em 1990 tinha 38.000.000 de habitantes, e em 2020 47.000.000
    \\[5pt]
    Para resolvermos, basta primeiro: Saber quanto a população desses países aumentou nesses intervalos de tempo. Descobrimos isso (obviamente) subtraindo a população no momento mais recente e no passado. No Brasil era 190.700.000 e aumentou para 211.800.000, então temos $211.800.000 - 190.700.000$. A população do Brasil aumentou 21.100.000 entre 2010 e 2020. Já a população da Espanha era 38.000.000 e aumentou para 47.000.000. Logo, calculamos $47.000.000-38.000.000$. Uma variação de 9.000.000
    \\[5pt]
    Feito isso, agora, basta calcular quanto o aumento de cada população representa de sua população total. Para isso, dividimos o aumento pela população, assim, descobrimos quantas unidades aumentaram para cada unidade da população. Começamos com o Brasil, que teve um aumento de 21.100.000 e tinha uma população de 190.700.000. Logo, calculamos $21.100.000/190.700.000$, simplificando, $211/1907$, que resulta em 0,11 aproximadamente. Já no caso da Espanha, tivemos um aumento de 9.000.000, e a população era de 38.000.000, logo, $9/38$, que resulta em aproximadamente 0,23. Então, pra cada unidade de poopulação do Brasil, tivemos 0,11 unidades de aumento, e para cada unidade de população da espanha, tivemos um aumento de 0,22. Ou seja, a Espanha teve um aumento percentual de população maior em relação ao Brasil
    \\[5pt]
    \textbf{Importante:} Tanto faz você calcular quanto o aumento representa da população antiga ou da mais recente. Ou seja, no caso do Brasil por exemplo, não importaria de calculássemos $21.100.000/211.800.00$ ou $21.100.000/190.700.000$, CONTANTO QUE, quando fossemos calcular quanto o aumento representa da população no caso da Espanha, também calculassemos em relação a popução da data que escolhemos no Brasil
    \\[5pt]
    \textbf{SEGUNDA FORMA DE RESOLVER O PROBLEMA:} Dividir a população final pela população inicial e analisar o fator de crescimento
    \[
    \text{Brasil: } 
    \frac{211{,}800{,}000}{190{,}700{,}000} \approx 1{,}1106 
    \;\Rightarrow\; 11{,}06\%.
    \]
    
    \[
    \text{Espanha: } 
    \frac{47{,}000{,}000}{38{,}000{,}000} \approx 1{,}2368 
    \;\Rightarrow\; 23{,}68\%.
    \]
    
    \[
    \therefore \text{A Espanha teve maior crescimento percentual.}
    \]
    \textit{Mas por que subtraímos 1 depois de calcular a razão de crescimento?}
    \\[5pt]
    Simples, Quando você divide a população final pela inicial: $\frac{P_i}{P_f}$ o resultado inclui o 100\% da população inicial + o crescimento
    \\[5pt]
    \textbf{Exemplo intuitivo:} Se algo cresce 10\%, então o valor final é:

    \[
    P_f = 1,10 \cdot P_i
    \]
    
    Note que 1,10 = 1 + 0,10.
    
    - O 1 representa os 100\% do valor inicial.
    - O 0,10 representa o crescimento de 10\%.
    
    Então, ao fazer:
    
    \[
    \frac{P_f}{P_i} = 1,10
    \]
    
    Você obtém:
    
    - 1 → a parte que já existia (100\% do início)
    - 0,10 → a parte que cresceu
    
    Por isso, para saber só o crescimento, você faz:
    
    \[
    \frac{P_f}{P_i} - 1
    \]

    
\end{itemize}

\subsection{Porcentagem}
Parte pelo todo. Depois, a fração da parte pelo todo, assim como (quase) toda fração, vai te dar um resultado entre [0, 1], então pra achar a porcentagem é só multiplicar por 100

\subsection{Divisão, Fração, Razão, proporção e suas aplicações}
Agora, chega de teoria! Você já sabe muito bem como a proporção funciona, com isso, você consegue identificar como aplicar ela! Vamos ver, na prática, como lidar com diferentes aplicações, formas rápidas de calcular problemas envolvendo esses assuntos, macetes, e tudo que envolve divisão

\begin{itemize}
    \item \textbf{Calcular Porcentagem}\\[5pt] Usar regra de 3 é burrice pra porcentagem em quase todos os casos. Sempre que pedir ``Ah, quantos \% tal quantidade representa de outra", é só usar parte pelo todo. A fração da parte pelo todo, assim como (quase) toda fração, vai te dar um resultado entre [0, 1], então pra achar a porcentagem é só multiplicar por 100
    \\[5pt]
    A não ser que: Peça coisas mais complexas. Aí tem que pensar mais

    \item \textbf{COMO CALCULAR PORCENTAGEM DE CABEÇA, DE FORMA RÁPIDA:}\\[5pt] Já pensou como você acha uma porcentagem arbitrária de forma fácil? Tipo, 27\% de 380. Simples! Você quebra esse desafio em partes. Pensa quanto seria 10\% de 380. Isso seria dividir 380 por 10, que é 38. Agora multiplica por 2, temos 76, sabemos quanto é 20\% de 380. Agora pense quanto é 5\% de 380. É a mesma coisa que a metade de 10\%, certo? Logo é 19. Soma isso com 72, temos 89. Isso representa 25\% de 380. Agora quanto é 1\% de de 380? Mesma coisa que dividir por 100, é 3,8. Agora multiplica isso por 2, obtemos 7,2 e somamos com o valor anterior, 89, obtendo 96,2. Pronto. Temos que 27\% de 380 é 96,2. Tá vendo, não é tão impossível 

    \item \textbf{Quando usar regra de 3 e quando usar ``Proporção por Lógica de Unidade"?}\\[5pt] Em geral, quando você quer calcular de cabeça, usa Proporção por Lógica de Unidade. E quando é no papel ou em casos que não daria pra fazer de cabeça, porque seria muito complicado, usa a regra de 3
\end{itemize}

\subsection{M.M.C, M.D.C e Aplicações}
Pô esse aí é importante. Da pra descobrir muitas coisas usando o conceito de MMC, tipo, quando 2 eventos vão acontecer simultaneamente
\\[5pt]
\textbf{Quando o MMC é a o produto entre os números?}
\\[5pt]
Apenas se os números (ou polinômios) forem primos entre si (ou seja, não compartilharem fatores primos). \textbf{Exemplo:} MMC entre 3 e 5 = 3 x 5 = 15, pois ambos os números são primos
\\[5pt]
\textbf{Mas por que é assim?}
\\[5pt]
Simples. Porque quando fazemos o MMC de números temos que fatora-los antes de multiplica-los, e no caso de número primos, não temos como fatorar. Logo, pra fazer o MMC entre eles é só multiplicar um pelo outro

\subsection{Propriedades e Macetes com Frações}

\textbf{Somar e Subtrair Frações}
\begin{itemize}
    \item USE O MÉTODO DA BORBOLETA. CASO CONTRÁRIO É MUITO DIFÍCIL
    \item\textbf{Método Formal (!!!!!!!!!!!!!!!! NÃO USE)}
    \\[5pt]
    Encontre faça o MMC entre os denominadores. As frações estando com o mesmo denominador, basta somar o numerador. Para encontrar o numeador, divida o MMC entre os denominadores pelo denominador de cada fração, e multiplique pelo numerador
    \\[10pt]
    \textbf{Por que não podemos somar frações com denominadores diferentes?}
    \\[5pt]
    Porque somar frações com denominadores diferentes é tipo tentar somar o valor da área de 2 pedaços de pizza, de 2 pizzas diferentes de mesmo tamaho, mas sem saber em quantos pedaços cada uma foi dividida. Primeiro, a gente tem que padronizar o tamanho de 1 pedaço, fazemos isso dividindo as pizzas na mesma quantidade de pedaços, garantindo que eles tem mesmo tamanho. Perceba que sem fazer isso, a gente não saberia a área de 1 pedaço de pizza, podia ser que a primeira estivesse dividida em 2 pedaços (área enorme) e a outra em 20 pedaços (área minúscula). Padronizando em quantos pedaços ambas as pizzas foram divididas, sabemos que a área de 1 pedaço da pizza A = área do pedaço da pizza B. Aí podemos somar porque sabemos que as áreas são iguais, teremos algo como x + x = 2x  
\end{itemize}

\textbf{Dividir é a mesma coisa que multiplicar pelo inverso}
\\[5pt]
Isso é \textbf{muito} útil na hora de manipular expressões, especialmente em cálculo 1. Olha esse exemplo:
\\[5pt]
\[
\frac{x^3}{\frac{x^2+1}{x}}=\frac{x^3}{x^2+1}\cdot\frac{1}{x}
\]
Nesse caso, a divisão é direta, ela já tá ali e você aplicou direto, nas em alguns casos, você tem que "caçar" uma divisão, e aí inverter
\\[10pt]
\textbf{Propriedade da subtração de frações com mesmo denominador}
\[
\frac{a - b}{c} = \frac{a}{c} - \frac{b}{c}
\]
Apesar do contrário ser óbvio, desse jeito você não costuma lembrar, então melhor deixar escrito
\\[10pt]
\textbf{Multiplicar e Dividir por um mesmo valor}
\[
\frac{A}{B} = \frac{A}{C} \cdot \frac{C}{B} \quad \text{(desde que } C \neq 0\text{)}
\]
É útil na hora de manipular expressões e polinômios
\\[10pt]

\vspace{15pt}
\hrule

\section{Macetes Gerais da Matemática}
\subsection{Somar frações com números inteiros de forma rápida}
Fácil! Basta multiplicar o inteiro pelo denominador e somar com o numerador. Isso será o numerador da fração resultante dessa soma, e o denominador você apenas repete\\[10pt]
Em termos matemáticos, temos que: 
\[
a + \frac{b}{c} = \frac{a\times c + b}{c}
\]
\textbf{Veja um exemplo:}\\[10pt]
$2 + \frac{8}{9}$ , multiplique o 2 (que é o inteiro) pelo 9 (denominador) e some com 8 (numerador)\\[10pt]
$2 + \frac{8}{9} = \frac{2 \times 9 + 8}{9} = \frac{26}{9}$ , após o passo anterior, apenas repetimos o denominador na fração resultante da soma. Simples assim

\subsection{Racionalização e conjugado}
\textbf{Racionalizar} é \textbf{tirar a raiz do denominador} (não pode deixar raiz embaixo da fração).

\begin{itemize}
  \item \textbf{Se o denominador é uma raiz só}, multiplica \textbf{em cima e embaixo pela mesma raiz}.
  \[
  \frac{1}{\sqrt{2}} \cdot \frac{\sqrt{2}}{\sqrt{2}} = \frac{\sqrt{2}}{2}
  \]

  \item \textbf{Se o denominador tem soma/subtração com raiz}, multiplica pelo \textbf{conjugado} (troca o sinal do meio).
  \[
  \frac{1}{2 + \sqrt{3}} \cdot \frac{2 - \sqrt{3}}{2 - \sqrt{3}} = \text{(vai sumir a raiz!)}
  \]
\end{itemize}
Se tiver mais de 2 termos no denominador, esse método não funciona :((((
\\[10pt]
\textbf{Racionalização para Raizes Cúbicas}
\\[5pt]
Pra raízes cúbicas, você tem que usar o fator racionalizante de raízes cúbicas. É a mesma ideia, só que mais difícil
\vspace{15pt}
\hrule

\section{Trigonometria}

\subsection{Trigonometria no Triângulo Retângulo}

\subsection{Trigonometria na Circunferência}

\subsection{Equações Trigonométricas}

\vspace{15pt}
\hrule

\section{Fatoração de Polinômios}
As fatorações de polinômios são uma forma de manipulação algébrica que permite simplificar ou reescrever polinômios como um produto de fatores menores. Isso facilita ao resolver equações, identificar raízes, analisar o comportamento do polinômio, realizar operações como divisão e integração de forma mais eficiente, dentre outras muitas utilidades.
\\[10pt]
\textbf{Temos diferentes casos de fatoração de polinômios, que serão apresentados a seguir}
\subsection{Diferença de Dois Quadrados}
Essa é muito fácil, é só isso:
\[
a^2 - b^2 = (a - b)(a + b)
\]
Sempre que tiver uma diferença entre dois números elevados ao quadrado, você pode aplicar essa propriedade, e quebrar a expressão em um produto de uma diferença por uma soma

\subsection{Trinômio Quadrado Perfeito (TQP)}
Quando você tem polinômios, mais especificamente, trinômios, que são polinômios com 3 termos, TALVEZ seja aplicável a fatoração por TQP. Esse caso de fatoração é aplicável somente se o trinômio tem a forma:
\[
a^2 \pm 2ab + b^2
\]
Ou seja, simplificando, se o trinômio satisfaz as seguintes condições:
\begin{itemize}
    \item O primeiro termo é um quadrado perfeito (\(a^2\)).
    \item O terceiro termo é um quadrado perfeito (\(b^2\)).
    \item O termo do meio é o dobro do produto das bases dos quadrados (\(\pm 2ab\)).
\end{itemize}
\textbf{Mas o que é um quadrado perfeito?}
\\[10pt]
Um \textbf{quadrado perfeito} é um caso especial onde o resultado da potência é um número inteiro ou uma expressão com raízes exatas. Em outras palavras, o valor do quadrado perfeito pode ser expresso como o quadrado exato de outro número ou expressão simples.
\\[10pt]
\textbf{Exemplos:}
\begin{itemize}
    \item \(4\) é um quadrado perfeito (\(2^2 = 4\)).
    \item \(25\) é um quadrado perfeito (\(5^2 = 25\)).
    \item \(x^2\) é um quadrado perfeito porque é o quadrado de \(x\).
\end{itemize}
Agora voltando ao tópico principal. A fatoração, se aplicável, terá a seguinte forma:
\[
a^2 \pm 2ab + b^2 = (a \pm b)^2
\]
\textbf{Exemplo:} \\
\[
x^2 + 6x + 9
\]
- \(x^2\) é um quadrado perfeito (\(x\)). \\
- \(9\) é um quadrado perfeito (\(3\)). \\
- \(6x = 2 \cdot x \cdot 3\). \\

\textit{Fatoração:} \\
\[
x^2 + 6x + 9 = (x + 3)^2
\]

\subsection{Fator Comum em Evidência}
A \textbf{fatoração por fator comum em evidência} é uma técnica utilizada para simplificar expressões algébricas onde os termos possuem um fator comum. O processo consiste em identificar o maior fator comum a todos os termos e colocá-lo em evidência, resultando em um produto entre o fator comum e o que resta da expressão original.

\subsection*{Passos:}
\begin{enumerate}
    \item \textbf{Identificar o fator comum}: Procure um fator que apareça em todos os termos da expressão.
    \item \textbf{Fatorar}: Extraia esse fator comum e escreva-o fora dos parênteses.
    \item \textbf{Escrever a expressão fatorada}: Dentro dos parênteses, coloque os termos que restaram depois de dividir cada termo original pelo fator comum.
\end{enumerate}

\subsection*{Exemplo:}
Dada a expressão \( 6x^2 + 9x \), o fator comum é \( 3x \). A fatoração será:

\[
6x^2 + 9x = 3x(2x + 3)
\]

\subsection{Fatoração por Agrupamento}
A fatoração por agrupamento é uma técnica útil quando se trata de polinômios com quatro ou mais termos. O objetivo é agrupar os termos de forma que possamos fatorar fatores comuns dentro de cada grupo. Aqui está um resumo do processo:

\begin{enumerate}
    \item \textbf{Agrupar os termos}: Divida os termos do polinômio em dois grupos. Por exemplo, considere o polinômio \( ax + ay + bx + by \).
    
    \item \textbf{Fatorar os grupos}: Em cada grupo, extraia o fator comum. No exemplo acima, teríamos:
    \[
    a(x + y) + b(x + y)
    \]
    
    \item \textbf{Fator comum}: Agora, observe que \( (x + y) \) é um fator comum entre os dois termos. Então, o polinômio pode ser fatorado como:
    \[
    (a + b)(x + y)
    \]
\end{enumerate}

Essa técnica é especialmente útil em expressões mais complexas, onde o agrupamento leva a fatores comuns que não são evidentes à primeira vista.

\subsection{Fatoração de Equações Quadráticas Complicadas}
A fórmula de Bhaskara (ou fórmula quadrática) é usada para resolver a equação quadrática. Ela nos dá as raízes da equação, que chamaremos de \(r_1\) e \(r_2\):

\[
x = \frac{-b \pm \sqrt{b^2 - 4ac}}{2a}
\]

Aqui:
\begin{itemize}
    \item \(b^2 - 4ac\) é chamado de discriminante (\(\Delta\)), e ele determina o número de soluções:
    \begin{itemize}
        \item Se \(\Delta > 0\), temos duas raízes reais e distintas.
        \item Se \(\Delta = 0\), temos uma única raiz real (duas raízes iguais).
        \item Se \(\Delta < 0\), não temos raízes reais (apenas complexas).
    \end{itemize}
\end{itemize}

\subsection*{2. Encontrando as raízes \(r_1\) e \(r_2\)}

Uma vez que tenhamos as raízes \(r_1\) e \(r_2\), podemos fatorar o polinômio quadrático na forma:

\[
ax^2 + bx + c = a(x - r_1)(x - r_2)
\]

Essa forma é a fatoração completa do polinômio.
\textbf{Conclusão:}

\begin{itemize}
    \item O método de \textbf{Bhaskara} é muito útil quando a fatoração simples (soma e produto) não é viável.

    \item Depois de encontrar as raízes \(r_1\) e \(r_2\) usando a fórmula quadrática:
    \[
    x = \frac{{-b \pm \sqrt{{b^2 - 4ac}}}}{2a}
    \]
    você pode escrever o polinômio fatorado na forma:
    \[
    a(x - r_1)(x - r_2)
    \]

    \item Esse método é \textbf{universal} e funciona para qualquer equação quadrática! Qualquer uma do tipo \(ax^2 + bx + c\).

    \textbf{Importante: Muitas vezes é mais rápido usar soma e produto para encontrar as raízes}
    
\end{itemize}

\subsection{Dispositivo de Briot-Ruffini}

\subsection{Conclusões Finais (IMPORTANTE LER ISSO)}
Então, no fim das contas, você muitas vezes, pra fatorar vai ter que fazer algumas manipulaçõe algébricas pra cair em algum caso de faotoração conhecido e/ou usar outros métodos de fatoração pra depois usar OUTRO método e aí fatorar de vez a expressão, tipo uma fatoração aninhada
\\[5pt]
Mas, em uma regra geral, da pra sempre usar um tal caso de fatoração sempre pra umas determinadas expressões:
\begin{itemize}
    \item \textbf{SE FOR UM CASO} $a^2-b^2$ \\[5pt]
    Use sempre \textbf{diferença de dois quadrados}. Lembra que as vezes
    \item \textbf{SE FOR UM CASO} $a^2\pm 2ab+b^2$ \\[5pt]
    Ou vai ser \textbf{Trinômio Quadrado Perfeito} ou aquele caso de \textbf{Fatoração por Bháskara}
    \item \textbf{SE FOR UM CASO BIZARRO QUE NÃO ENCAIXA EM NENHUM CASO, E VOCÊ ESTIVER QUASE PEDINDO SOCORRO} \\[5pt]
    Primeiro, calma. Provavelmente você tem que fazer alguma manipulação algébrica pra chegar em algum dos casos conhecidos. Vê se usar coisas como \textbf{racionalização}, \textbf{multiplicar em cima e baixo da fração por um mesmo valor}, \textbf{substituir parte da expressão por} $u$ não resolve. Caso não, provavelmente não tem que fatorar mesmo e você estava indo pelo caminho errado, ou sei lá, mas saiba que isso é muito raro. É isso beijos XOXO
\end{itemize}
No fim, o que importa não é fatorar ou desenvolver uma expressão, é \textbf{simplificar.} Perceba, o conceito de simplificar não é equivalente a fatorar, fatorar simplifica? Sim, muitas vezes, afinal você está encurtando a expressão, mas muitas vezes, pra por exemplo, cancelar 2 termos de uma fração, vale mais a pena desenvolver uma expressão. Entendeu? O conceito de simplificar é: \textbf{Olhe, e veja o que é mais estratégico}. Não existe uma regra geral, cada caso é um caso. Mas é sempre importante lembrar, simpplificar \textbf{não é necessariamente} encurtar uma expressão, é ver uma estratégia para tornar tal problema mais fácil de resolver
\vspace{15pt}
\hrule
\vspace{15pt}
\section{Sistemas de Equação}
É um conjunto de equações. Se uma equação tem 1 única expressão e 2 variáveis, fica impossível determinar o valor numérico das 2. É possível determinar certas coisas em função de x ou y, mas determinar seus valores, não
\\[5pt]
A solução de um sistema é a solução comum para todas as suas equações
\\[5pt]
Para podermos determinar os valores das variáveis, temos que ter 1 equação para cada variável, formando um sistema
\subsection{Sistemas 2x2}
Usa o método da substituição ou adição. Mais raramente, usa-se o método da comparação
\\[5pt]
A solução de um sistema 2x2 é sempre um \textbf{par ordenado}

\subsection{Sistemas 3x3}
Sim, é possível resolver sistemas 3x3 utilizando o método da substituição, mas \textbf{não é nem um pouco prático}
\\[5pt]
A solução de um sistema 3x3 é sempre uma \textbf{tripla ordenada}
\\[5pt]
Temos vários métodos para resolver sistemas 3x3 (e maiores). Vamos apresenta-los:
\begin{itemize}
    \item \textbf{Método do Escalonamento:}
    \\[5pt]
    É provavelmente o melhor método para resolver sistemas 3x3 (e maiores até!)
    \\[5pt]
    Primeiramente, veja como é fácil resolver um sistema "escada"
    \[
\begin{cases}
x + 2y + z = 4 \\
\phantom{x + {}} y + z = 1 \\
\phantom{x + {} y + {}} 2z = 6
\end{cases}
\]
Para encontrar $z$
\[
2z = 6 \Rightarrow z = 3
\]
Para encontrar $y$
\[
y + 3 = 1 \Rightarrow y = -2
\]
Para encontrar $x$
\[
x + 2(-2) + 3 = 4 \Rightarrow x - 4 + 3 = 4 \Rightarrow x = 5
\]
\[
(x,\ y,\ z) = (5,\ -2,\ 3)
\]
Viu só? É só resolver da equação mais simples pra mais complexa. No caso, primeiro encontramos $z$ e a partir do valor de $z$ determinamos o valor de todas as outras variáveis
\\[5pt]
A ideia do escalonamento é fazer o sistema virar uma "escada", através da soma de equações de forma que as variáveis se anulem, ou seja, criar um sistema escada, fácil de resolver, igual o do exemplo anterior. Essas são as etapas do escalonamento:
    \begin{itemize}
        \item Primeiro passo: Fixe uma equação. De preferência a equação mais simples do sistema, e que tenha $x$. Simples assim. Copie ela
    \end{itemize}
\end{itemize}

\section{Equações do Segundo Grau}
São equações que tem pelo menos uma incógnita elevada ao quadrado
\subsection{Resolução pelo Método da Soma e Produto}
É bem simples, veja na equação quanto vale $[a, b, c]$ e calcule os seguintes valores:
 \[
    S = -\frac{b}{a} \quad P = \frac{c}{a}
\]
Após calcular os valores de S e P, é só pensar: "Quais valores quando somados resultam no valor de S e quando mutiplicados resultam no valor de P?"
\\
Exatamente por isso que o valor S é chamado de valor da "Soma" e o valor P de "Produto". \textbf{DICA: Sempre pense primeiro em quais valores quando multiplicados resultam em P, costuma facilitar o raciocíneo}


\section{Somatório}
\[
\sum_{i=1}^{n} x_i = x_1 + x_2 + \dots + x_n
\]
Esse é um somatório de $x$ índice $i$, ou seja, você vai somar todos os índices de o índice $i$, que no caso é 1, até n
\\[5pt]
Ó, pra ler um somatório é fácil. O número em baixo é onde ele começa, o número de cima é até onde ele vai, e o número (ou expressão) na frente dele indexa o valor que está sendo somado a uma variável, e faz alguma coisa. \textbf{Importante:} Você tem \textbf{sempre} que indexar o valor de baixo (no caso o $i$), porque se não seu somatório não faz nada
\\[5pt]
Se esse somatório fosse um loop \textit{for}, seria assim: \textit{for(i=1, $i < n$, i++)}

\section{Inequações}
Inequação é uma sentença matemática, com 1 ou mais incógnitas, e que contenha uma desigualdade, diferente das equações, que representam uma igualdades. Elas representam relações que não são de equivalência. \textbf{Ou seja, inequações, diferente das equações, não te fornecem um valor, mas sim, um intervalo de valores}

\subsection{Inequações do 1º grau}
Você resolve que nem uma equação do 1º grau. Só tem uma diferença, quando você multiplicar ambos os membros da inequação por valores negativos, você deve inverter o sinal da desigualdade

\subsection{Inequações Produto e Inequações Racionais}
São as inequações que tem polinômios contendo produtos e frações

\subsection{Inequações do 2º grau}
\begin{itemize}
    \item EM CASO DE POLINÔMIOS DO 2º GRAU SIMPLES: Para obter o conjunto solução desse tipo de inequação, extraia as raízes e esboce o gráfico. Ao realizar essas duas etapas você fez o que se chama de \textbf{estudo do sinal}
    \item  A função principal pode ser muito complexa, e envolver polinômios que poderiam ser separados como outras funções polinomiais, exemplo: \[
\frac{(-x + 3)(x^2 - 3x + 2)}{2x - 1} < 0
\] em casos assim, você primeiramente coloca esse polinômio em função de alguma variável, x nesse caso, e decompõe em funções menores, criando as funções: \[
f(x) = -x + 3
\]
\[
g(x) = x^2 - 3x + 2
\]
\[
h(x) = 2x - 1
\] após separar a função em funções menores, estude o sinal de cada uma dessas funções individualmente, e então, crie uma tabela com os sinais de cada função em cada intervalo, e multiplique todos os sinais para obter o estudo do sinal da função principal, ficaria assim
\[
\begin{array}{|c|c|c|c|c|c|}
\hline
x & \frac{1}{2} & 1 & 2 & 3 \\ \hline
f(x) & + & + & + & - \\ \hline
g(x) & + & + & - & + \\ \hline
h(x) & - & + & + & + \\ \hline
\frac{f(x) \cdot g(x)}{h(x)} & - & + & - & - \\ \hline
\end{array}
\]
\end{itemize}

\section{Funções}

\subsection{Noção Intuitiva de Função}
De forma simples, uma função é uma regra que relaciona cada elemento de um conjunto (chamado domínio) a exatamente um elemento de outro conjunto (chamado contradomínio). É uma regra que relaciona elementos de 2 conjuntos
\\[5pt]
Vamos dizer que o domínio é um conjunto $A$ e o contradomínio um conjunto $B$, nesse caso, escrevemos a função como:
\[
f: A \to B
\]
Essa notação mostra que $f$ pega elementos de $A$ e transforma em elementos de $B$
\\[5pt]
Na notação de função, chamamos o elemento (valor) de entrada de $x$, e o valor de saída (correspondente a $x$) de $f(x)$, ou $y$.
\\[10pt]
\subsection{Notações de Função}

\textbf{Mas por que 2 notações para o valor de saída?}
\\[5pt]
Porque cada uma é útil em um contexto. A notação $f(x)$ é útil quando queremos explicitar qual foi o valor de entrada $(x)$ da função. Ex: $f(1), f(a), f(288)$. Isso é útil quando queremos enfatizar a relação entre entrada e saída, como em cálculos e manipulações algébricas. Já a notação $y$ é mais resumida sem precisar mencionar explicitamente a entrada $x$. Isso é útil em gráficos e equações onde o foco está mais na saída do que na entrada
\\[5pt]
Ainda nesse tópico, funções podem ter muitas notações. Por exemplo:
\begin{itemize}
    
    \item \textbf{Notação de Diagrama de Setas:} A dos círculos com setinhas, usada para representar o domínio, contradomínio e a regra de associação visualmente. Mostra dois conjuntos com setas ligando os elementos do domínio aos do contradomínio.
    
    \item \textbf{Notação por Regras:} A mais comum na matemática formal. Escrevemos algo como \( f(x) = 2x + 1 \), indicando como calcular a imagem de cada elemento.
    
    \item \textbf{Notação de Definição Formal:} Indica domínio, contradomínio e a regra ao mesmo tempo. Exemplo: 
    \[
    f: \mathbb{R} \to \mathbb{R}, \quad f(x) = x^2
    \]
    
    \item \textbf{Notação de Pares Ordenados:} Representa a função como um conjunto de pares \( (x, f(x)) \). Exemplo: 
    \[
    f = \{ (1, 2), (2, 4), (3, 6) \}
    \]
\end{itemize}

A notação mais comum para função é:
\[
f: A \to B
\]
Isso indica que a função $f$ leva um elemento $x \in A$ a $f(x) \in B$

\subsection{Definição Formal de Função}
Uma função é composta por uma terna, $(f, A, B)$, onde $f$ é a lei que define a função, chamada \textbf{Lei de Formação}, $A$ é o domínio (valores de entrada) de $f$ e $B$ o contra-domínio (valores de saída). A lei $f$ associa a cada $x$ um \textbf{único} $f(x)$ (podemos chamar $f(x)$ de $y$ também) pertencente a $B$
\\[10pt]
Dito isso, definimos função assim:
\\[10pt]
\textbf{Seja \( A \) um conjunto, chamado domínio, e \( B \) outro conjunto, chamado contradomínio. Dizemos que \( f : A \to B \) é uma função se todo elemento de \( A \) está associado a um único elemento de \( B \)}
\\[10pt]
Observe que, na notação de função que usamos, $f$ representa a função mesmo, a relação em si, $f(x)$ representa os valores do contra-domínio e $x$ os valores do domínio

\subsection{Domínio, Contradomínio e Imagem}
Quando estamos tratando de funções, damos alguns nomes específicos para os conjuntos que a \textbf{lei de formação} da função relaciona
\begin{itemize}
    \item \textbf{Domínio ($\mathrm{D}$)}
    \\[5pt]
    \textit{O domínio é conjunto de todos os valores de ``entrada" da função. É o conjunto de valores para os quais a função está definida}
    \\[5pt]
    Calma. Explicando melhor, as vezes o domínio é um tal conjunto $C$ porque a função apenas está definida em $C$, logo o domínio é restrito a $C$ por esse motivo. No entanto, em muitos casos, a função está definida em $\mathbb{R}$, porém queremos estudar apenas os valores dela nos números inteiros, por exemplo. Nesse caso, por nossa vontade restringimos o domínio. E podemos restringir o domínio a qualquer conjunto arbitrário que quisermos, essa é a mágica. Se quisermos, podemos falar que o domínio é só um conjunto $H = \{0.0032, 1, 42, 877\}$ olha que conjunto nada a ver, com 4 elementos só, e ainda assim ele pode ser o domínio sem problemas
    \\[5pt]
    Se não for especificado o domínio de uma função, assume-se que o domínio é todo o conjunto dos números reais
    \\[5pt]
    \textbf{Exemplos:}
    \[f(x) = \frac{1}{x}\]
    Essa função não está definida para $x = 0$, logo, seu domínio é $\mathrm{Dom}(f):\mathbb{R} - \{0\}$ (todos os reais sem ser o 1)
    \[f(x) = 0,5 + x\]
    \[f:\mathbb{Z}\xrightarrow{}\mathbb{R}\]
    Agora, nesse exemplo, a função está definida para todo o conjunto $\mathbb{R}$, mas, por livre e espontânea vontade resolvemos restringir o domínio aos inteiros, e tudo certo

    \item \textbf{Contra-Domínio ($\mathrm{CD}$)}
    \\[5pt]
    É o conjunto de saída da função. Ele é definido pelo criador da função, junto com domínio. Você que escolhe o contra-domínio
    \\[5pt]
    Perceba! O conjunto de saída pode não ser \textit{somente} o conjunto de valores gerados pela função (essa é a imagem). O contra-domínio contém o conjunto imagem, essa é a única coisa que podemos afirmar, olha só
    \[
    \text{imagem} \subseteq \text{contradomínio}
    \]
    Ou seja, todo valor da imagem está dentro do contradomínio, mas nem todo valor do contradomínio precisa estar na imagem (a menos que a função seja sobrejetora)

    \item \textbf{Imagem ($\mathrm{Im}$)}
    \\[5pt]
    É o conjunto de todos os valores gerados pela função. Esse conjunto inclui todos os todos os valores $f(x)$ quando $x$ percorre todo o domínio
    \\[5pt]
    \textbf{Então, o contra domínio é a mesma coisa que a imagem? Não. Não é}
    \\[5pt]
    O contradomínio é o conjunto de chegada da função. É o conjunto no qual os valores da função estão contidos, ou seja, onde a função "despeja" os resultados, não necessariamente \textbf{somente} os resultados. Isso é, o contra domínio pode ter mais elementos que a imagem gerada pela função \textbf{Importante:} O contradomínio é definido junto com a função, não é descoberto depois. Ele não depende dos valores que realmente saem da função — isso é outra coisa, essa coisa é justamente a imagem. A imagem é o conjunto de valores efetivamente produzidos pela função. Entendeu a diferença? Se a gente tivesse uma função que enfia farinha, e ela faz biscoitos, é como se o domínio fosse a farinha, o contra-domínio fosse a "bandeja" onde a função despeja os biscoitos, e a imagem os biscoitos em si
    \\[5pt]
    \textbf{Exemplo:}
    \[
    f: \mathbb{R} \to \mathbb{R}, \quad f(x) = x^2
    \]
    \begin{itemize}
        \item \textbf{Domínio:} $\mathbb{R}$
        \item \textbf{Contradomínio:} $\mathbb{R}$
        \item \textbf{Imagem:} $[0, \infty[$, porque $x^2$ nunca da um número negativo
    \end{itemize}
    \end{itemize}

\subsection{Representação de Tabela das Funções}
Você pode representar funções usando tabelas que relacionam os valores de x e y, ou seja, montar tabelas que relacionam valores do domínio ao contra domínio, OU SEJA, podemos montar uma tabela que relaciona elementos de 2 conjuntos, pronto. É isso.

\subsection{Significado de ``estar em função" na matemática}

Quando dizemos, na matemática, que tal número ou expressão está em função de outra, queremos dizer que há (ou que definimos) uma relação de dependência hierárquica entre elas. Ou seja, se $Y$ está em função de $X$, isso significa que $X$ ``manda'' em $Y$ — o valor de $Y$ depende do valor de $X$.
\\[5pt]
Quando usamos expressões como ``resolver em função de tal variável'', estamos dizendo que queremos expressar todas as outras variáveis em termos daquela, \textbf{ela é o que importa}. Por exemplo, ao resolver ``em função de $x$'', significa que $x$ é a variável principal e queremos isolá-la ou entender como as outras dependem dela.

\subsection{Funções como Encapsuladores}
Sim, é basicamente isso que você leu. Muitas vezes, funções são usadas unicamente como "encapsuladores", principalmente em cálculo
\\[5pt]
\textbf{Mas o que é usar uma função como encapsulador?} Isso acontece principalmente quando queremos "guardar" um polinômio, expressão, ou até mesmo outra função, dentro de um objeto. Objeto esse que é, justamente, uma função. Em resumo, se você quiser condensar uma expressão em um objeto, para estudar seu comportamento, ou para simplificar/manipular expressões, ou para poder fazer composições de funções complexas, ou simplesmente para ocupar menos espaço, você usa funções
\\[5pt]
\textbf{Exemplos:}
\begin{itemize}
    \item \textbf{Encapsulando uma expressão algébrica:} \\
    Suponha a expressão:
    \[
    \frac{x^2 - 1}{x + 3} + \ln(x)
    \]
    Em vez de escrevê-la toda vez, definimos uma função \( f \) que a encapsula:
    \[
    f(x) = \frac{x^2 - 1}{x + 3} + \ln(x)
    \]
    \textbf{Utilidade:}
    \begin{itemize}
        \item Avaliar \( f(2) \) diretamente.
        \item Facilitar operações como derivar (\( f'(x) \)) ou integrar (\( \int f(x) \, dx \)).
    \end{itemize}

    \item \textbf{Funções para composição:} \\
    Para calcular \( \sin(2x + \cos(x)) \), podemos definir:
    \[
    g(x) = 2x + \cos(x) \quad \text{e} \quad h(x) = \sin(g(x))
    \]
    \textbf{Vantagem:}
    \begin{itemize}
        \item Separar a operação em partes gerenciáveis.
        \item Aplicar a regra da cadeia na derivação: \( h'(x) = \cos(g(x)) \cdot g'(x) \).
    \end{itemize}

    \item \textbf{Polinômios como funções:} \\
    Considere o polinômio:
    \[
    P(x) = x^3 - 4x^2 + 2x - 1
    \]
    \textbf{Aplicações:}
    \begin{itemize}
        \item Encontrar raízes (\( P(x) = 0 \)).
        \item Estudar comportamento assintótico sem reescrever o polinômio.
    \end{itemize}

    \item \textbf{Funções definidas por partes:} \\
    Exemplo:
    \[
    f(x) = 
    \begin{cases} 
    x^2 & \text{se } x \geq 0, \\
    -2x & \text{se } x < 0.
    \end{cases}
    \]
    \textbf{Por que encapsular?}
    \begin{itemize}
        \item Unificar regras em um único objeto (\( f \)).
        \item Simplificar análises de continuidade ou diferenciabilidade.
    \end{itemize}

    \item \textbf{Funções de múltiplas variáveis:} \\
    Em cálculo multivariável:
    \[
    F(x, y) = x^2 y + e^{xy}
    \]
    \textbf{Utilidade:}
    \begin{itemize}
        \item Calcular derivadas parciais (\( \frac{\partial F}{\partial x} \)).
        \item Avaliar \( F(1, 0) \) sem substituições manuais.
    \end{itemize}

    \item \textbf{Importância do encapsulamento:}
    \begin{itemize}
        \item Economia de notação (evita repetição).
        \item Abstração de operações complexas (ex.: \( f \circ g \)).
        \item Clareza em provas e manipulações.
    \end{itemize}
\end{itemize}

\subsection{Função Linear x Função Afim}
\begin{itemize}
    \item \textbf{Função afim}
    \\[5pt]
    Toda função da forma $f(x) = ax + b$ é dita afim e seu gráfico é uma reta
    \item \textbf{Função Linear}
    \\[5pt]
    Dada $f(x) = ax$ é dita função linear e seu gráfico é uma reta passando por (0, 0), ou seja, a origem do plano de ordenadas
\end{itemize}

\subsection{Funções Monótonas}
Funções monótonas são funções estritamente crescentes ou decrescentes, ou nem crescentes nem decrescentes (constantes). Elas são previsíveis, por isso ``monótonas"
\\[10pt]
Uma função $F: I \subset R \xrightarrow{} B \subset R$ é dita estritamente crescente em $I$ se para todo $x, y \in I$ temos que se $x < y$ então $f(x) < f(y)$. Ou seja, quanto maior o valor que você enfia, maior o valor que sai
\\[10pt]
Esta propriedade se reflete no gráfico de $F$

\begin{itemize}
    \item Uma função $F: I \subset R \xrightarrow{} R$ é dita estritamente decrescente em I se para todo $x, y \in I$ tal que $x < y$ temos $f(x) < f(y)$
    
    \item Uma função $F: I \subset \mathbb{R} \to B \subset \mathbb{R}$ é dita não crescente (ou decrescente) em $I$ se, para todo $x, y \in I$, temos que se $x < y$, então $F(x) \geq F(y)$.
\end{itemize}

Analogamente definimos função não crescente como sendo aquelas satisfazendo a seguinte propriedade
\\[5pt]
$para todox, y \in I = dom(F)$ temos que se $x < y$ temos que $f(x) maiorigual f(y)$

\subsection{Gráficos e Curvas}
\textbf{``Será que, dada uma curva no plano xy, sempre existe uma função cujo gráfico é aquela curva?"}
\\[5pt]
\textbf{Resposta:} Não. Porque podemos ter curvas que associam um elemento a mais de uma imagem, um x a mais de um y. Formas como um círculo no gráfico fazem isso
\\[10pt]
\textbf{Como "navegar" pelo gráfico de uma função?}
\\[5pt]
Basicamente, pra verificar qual valor a função retorna com um determinado valor, ou quanto a função vale em certo ponto (geralmente o estamos falando de um ponto nas abscissas), basta encontrar aquele ponto no eixo x, e ver onde ele "bate" no eixo y, porque lembre-se, o eixo x são os valores que estamos inserindo na função, e o eixo y os valores que a função está retornando 

\subsection{Tipos de Funções}
\begin{itemize}
    \item \textbf{Função sobrejetora:}\\[10pt] Uma função é sobrejetora se, e somente se, o seu conjunto imagem for especificadamente igual ao contradomínio, Im = B. Por exemplo, se temos uma função f : Z→Z definida por y = x +1 ela é sobrejetora, pois Im = Z. Ou seja, não ficou ninguém "de fora" no contra domínio

    \item \textbf{Função Injetora}\\[10pt] Uma função é injetora se os elementos distintos do domínio tiverem imagens distintas. Por exemplo, dada a função f : A→B, tal que f(x) = 3x.
    \\[5pt]
    Formalizando, dada $f: I \subset R \xrightarrow{} R$ dizemos que $f$ é injetora se para todo $x_1, x_2 \in I$ tal que $f(x_1) = f(x_2)$ então $x_1 = x_2$

    \item \textbf{Função bijetora:}\\[10pt] Uma função é bijetora se ela é injetora e sobrejetora. Por exemplo, a função f : A→B, tal que f(x) = 5x + 4.
\end{itemize}

\subsection{Classificação das Funções}
\begin{itemize}
    \item \textbf{Função Afim} \\[10pt] Você sabe o que é, você sabe
    \item \textbf{Função Polinomial} \\[10pt] É a função do tipo $y = ax^n + ax^n-1 + a_n-1$
    \item \textbf{Função Racional} \\[10pt] É o quociente de 2 funções afim
    \item \textbf{Função Trigonométrica} \\[10pt] Tem alguma coisa de trigonometria na lei de formação
    \item \textbf{Função Logaritmo} \\[10pt] Tem log na lei de formação 
    \item \textbf{Função Exponencial} \\[10pt] É a inversa da logaritmo
\end{itemize}

\subsection{Função Composta}
A função composta acontece quando o resultado de uma função serve como entrada para outra.  
Seja \( f: B \to C \) e \( g: A \to B \), então a composição \( f \circ g \) é uma função de \( A \) em \( C \), definida por:
\[
(f \circ g)(x) = f(g(x)) \quad \text{para todo } x \in A
\]
\textbf{Exemplo:}  
Se \( f(x) = x^2 \) e \( g(x) = x + 1 \), então:
\[
(f \circ g)(x) = f(g(x)) = f(x + 1) = (x + 1)^2
\]
\textbf{Dica para lembrar:}  
Na composição \( f \circ g \), \textbf{g age primeiro}, \textbf{f depois}
\\[10pt]
Para a notação de função composta calculada em um dado valor ou ponto, temos \( (f \circ g) (x)\). Essa notação diz "calcule $f$ composta em $g$ no valor (ou ponto) $x$". Primeiro o valor $x$ passa por $g$, depois por $f$

\subsection{Função Inversa}
A função inversa é um tipo de função que joga os elementos do contra-domínio no domínio, assim, invertendo a função original. Pensa assim, uma função ``normal" sempre liga um $x$ a um $y$. O que a função inversa faz é ligar cada $y$ a um $x$. Ou seja, a função inversa de $f$, denotada por $f^-1$ ``desfaz" o que $f$ faz\\[10pt]
Para que uma função seja inversível, ela deve ser \textbf{bijetora}. Vamos apresentar a definição formal de função inversa\\[10pt]
Uma função $f$ bijetora tem por inversa a função $f^{-1}$, tal que $(x, y) \in f \Longleftrightarrow (y, x) \in f^{-1}$\\[10pt]
\textbf{Explicando a definição:} O par $(x, y)$ na função $f$ contém $x$, que representa um elemento qualquer do domínio, e $y$ como a imagem correspondente a $x$ no contra domínio. Dito isso, o que a definição nos diz é: Se você tem esse par $(x, y)$ na função original, a função inversa irá inverter o elemento de entrada (elemento do domínio), que era $x$, com o elemento de saída (elemento do contra-domínio), que era $y$, ficando com um par $(y, x)$. Ou seja, a função $f^{-1}$ inverte os papéis do domínio e contradomínio da função original\\[10pt]
\textbf{Por que $f$ deve obrigatoriamente ser bijetora para ser inversível?}\\
Lembrando, função bijetora é aquela que é ao mesmo tempo sobrejetora e injetora. Isso é uma condição obrigatória para uma função ser inversível pois se 2 ou mais elementos do domínio se ligassem a um único elemento do contra-domínio, a função inversa de $f$ \textbf{não seria função!} Lembre-se da definição de função, "um elemento do domínio deve se ligar a \textbf{somente} um único elemento do contra domínio", caso contrário, não existe função\\[10pt]
Para calcular a função inversa, $f^{-1}$, em uma função $y = f(x)$, troca-se $x$ com $y$ e isola-se $y$\\[10pt]
Vamos ver alguns \textbf{exercícios} para entender melhor\\[10pt]
a) $y = 2x - 3$\\[5pt]Vamos trocar todos os $x$ por $y$\\[5pt]
$x = 2y - 3$\\[5pt]Agora que trocamos, vamos isolar $y$\\[5pt]
$-2y = -x -3$\\[5pt]$2y = x + 3$\\[5pt]$y = \frac{x + 3}{2}$\\[5pt]Isolamos $y$. Perceba que multiplicamos a equação toda por -1 para simplificar, e depois apenas resolvemos a equação\\[5pt]
Com isso, temos que a função inversa de $y = 2x - 3$ é $y = \frac{x + 3}{2}$. Na notação de função inversa, $f^{-1}(x) = \frac{x + 3}{2}$\\[20pt]
b) $f(x) = \frac{x + 4}{2}$\\[5pt] Lembre-se que $f(x)$ e $y$ são a mesma coisa, logo, podemos trocar a notação $f(x)$ por $y$

\subsection{Funções e o Plano Cartesiano}
Uma forma de representar um relação de função é através do plano cartesiano, onde um eixo representa os valores de entrada, e o outro os valores de saída

\subsection{Determinar se um Gráfico é Função}
É só fazer o famoso \textbf{Teste da Reta Vertical}. Ele funciona assim, você desenha uma reta vertical, perpendicular ao eixo das abscissas (x), e verifica se a reta cruza o gráfico da função em mais de 1 ponto. Se cruzar em mais de 1 ponto, não é função, porque perceba que se isso acontece, quer dizer que há mais de 1 valor $x$ para cada $y$, e isso fura a definição de função (cada elemento do domínio deve ter apenas 1 elemento correspondente no contra-domínio)

\subsection{Exercícios Úteis}
Uma certa \( f : \mathbb{R} \to \mathbb{R} \) tem a seguinte propriedade:
\[
f(mx) = mf(x) \quad \text{para todos } m, x \in \mathbb{R}.
\]
Descreva a expressão analítica de \( f(x) \).
\\[5pt]
\textbf{Resposta:} A ideia pra resolver é: A função tem que ser do tipo em que multiplicar o x antes ou multiplicar a saída depois dá o mesmo resultado. Ou seja: Se eu multiplico $x$ por $m$ antes de jogar na função, dá a mesma coisa que jogar $x$ na função e multiplicar o resultado por $m$ depois.
\\[5pt]
Então, como resposta dessa equação funcional \( f(x) = kx \) serve, porque se você colocar \( mx \) no lugar de \( x \), ela vira \( f(mx) = k(mx) = m(kx) = mf(x) \). Então a equação fica certa. Por isso, a resposta é \( f(x) = kx \), com \( k \in \mathbb{R} \).

\section{Logaritmos}

\subsection{Alongamentos Úteis}

\section{Progressão Aritimética e Geométrica}
\hrule
\vspace{10pt}
\section{Potências}
\subsection{Expoentes Negativos}
Inverte a base. Só.

\subsection{Expoentes Fracionários}
``O que tava no sol vai pra sombra, e o que tava na sombra vai pro sol"

\subsection{Expoentes Fracionários Negativos}
Similar ao caso anterior. Inverte a base, mas transforma em, raíz já que é um expoente fracionário. Vale notar que você tem que verificar se a fração é realmente negativa, exemplo $10^-4/-8$ \textbf{não é um expoente fracionário negativo}, afinal, menos com menos da mais. Mas em casos como $3^-2/5$, essa fração é realmente negativa, então invertemos
\vspace{10pt}
\hrule
\vspace{10pt}
\section{Raízes}
\subsection{Subtração de Raizes}
Para subtrair raizes, em quase todos os casos sua intençaõ é fatorar os números e resolver uma parte da raiz, ou seja, no caso de subtração de raizes sua intenção é simplificar a expressão ao máximo, e (geralmente) não exatamente resolver encontrando como resultado um número
\\[10pt]
Passos comuns para subtrair raízes:
\begin{itemize}
    \item Fatorar os números dentro da raiz para ver se há algo que pode ser simplificado.
    \item Extrair fatores quadrados (ou cúbicos, dependendo da raiz) para fora da raiz.
    \item Verificar se há termos semelhantes para realizar a subtração.
\end{itemize}
\textbf{Exemplo:}
\\[10pt]
$\sqrt{200} - \sqrt{32}$
\\[5pt] Pra resolver isso, vamos fatorar o 200 e o 32 pra ver se encontramos algum fator comum ou algo do tipo
\\[5pt]
$\sqrt{100} \cdot \sqrt{2} - \sqrt{16} \cdot \sqrt{2}$
\\[5pt]
Perceba que a fatoração fez aparecer raizes simples de ser resolver. Vamos resolve-las
\\[5pt]
$10 \cdot \sqrt{2} - 4 \cdot \sqrt{2}$
\\[5pt]
Agora, pra resolver essa subtração, vamos aplicar a mesma lógica que a aplicaríamos pra subtrações de variáveis
\\[5pt]
Imagine que ao invés de $\sqrt{2}$ fosse um $x$ ou $y$, por exemplo. Você saberia resolver, $10x - 4x = 6x$. Logo, pense como se $\sqrt{2} = x$, e resolva
\\[5pt]
$10 \cdot \sqrt{2} - 4 \cdot \sqrt{2} \xrightarrow{} 6\sqrt{2}$

\section{Conjuntos}
\subsection{Operações Entre Conjuntos}
\begin{enumerate}
    \item \textbf{União $(A \cup B)$}
    \\[5pt]
    Gera um novo conjunto com todos os elementos de A e todos os elementos de B
    \item \textbf{Interseção $(A \cap B)$}
    \\[5pt]
    Gera um novo conjunto que contém apenas os elementos que estão nos dois conjuntos ao mesmo tempo
    \item \textbf{Diferença $(A - B)$}
    \\[5pt]
    Gera um novo conjunto que contém os elementos que estão em A, mas não em B
\end{enumerate}
\textbf{Dica: O Diagrama de Venn ajuda a compreender essas operações e suas utilidades}
\begin{figure}[h]
    \centering
    \includegraphics[width=0.5\linewidth]{venn.jpg}
    \caption{Operação com os conjuntos}
    \label{fig:enter-label}
\end{figure}

\section{Coisas Envolvendo Divisão (Razão, Proporção, Porcentagem e mais)}
Agora: Temos um alimento de 250g com 45g de fibras, e outro de 100g com 20,5g. Qual deles tem mais fibras por grama?
Resolução: Divida a parte pelo todo pra saber quanto a parte representa do todo, e depois compare-as. A divisão que resultar em um resultado maior será o alimento com mais fibras. Parte dividida pelo todo no caso seria 45 ÷ 250 e 20,5 ÷ 100
\\[10pt]
\textbf{Por que essa estratégia funciona?} Porque obtivemos a quantidade de fibras por grama de cada alimento, e depois comparamos os dois pra ver qual tem mais fibras a cada grama, e óbvio, o tem mais fibras a cada grama tem mais fibras no total
\\[10pt]
\textbf{A ideia das razões é sempre: "Quantas unidades desse tem pra cada unidade daquele?"}

\section{Matrizes}

\subsection{Matrizes}
São tabelas com valores

\subsection{Matriz transposta, simétrica e antissimétrica}
\begin{itemize}
    \item \textbf{Matriz Transposta}:
    \begin{itemize}
        \item \textbf{Definição}: Troca as linhas pelas colunas.
        \item Se a matriz original é \( A \), a transposta é \( A^T \).
        \item Exemplo:
        \[
        A =
        \begin{bmatrix}
        1 & 2 \\
        3 & 4
        \end{bmatrix}
        \Rightarrow
        A^T =
        \begin{bmatrix}
        1 & 3 \\
        2 & 4
        \end{bmatrix}
        \]
    \end{itemize}

    \item \textbf{Matriz Simétrica}:
    \begin{itemize}
        \item \textbf{Definição}: É uma matriz que é igual à sua transposta: \( A = A^T \).
        \item Só existem matrizes simétricas \textbf{quadradas} (mesmo número de linhas e colunas).
        \item Exemplo:
        \[
        A =
        \begin{bmatrix}
        1 & 2 \\
        2 & 3
        \end{bmatrix}
        \Rightarrow A = A^T
        \]
    \end{itemize}

    \item \textbf{Matriz Antissimétrica}:
    \begin{itemize}
        \item \textbf{Definição}: A transposta é o oposto da matriz: \( A^T = -A \).
        \item Os elementos da diagonal principal são sempre \textbf{zero}.
        \item Exemplo:
        \[
        A =
        \begin{bmatrix}
        0 & 5 \\
        -5 & 0
        \end{bmatrix}
        \Rightarrow A^T =
        \begin{bmatrix}
        0 & -5 \\
        5 & 0
        \end{bmatrix} = -A
        \]
    \end{itemize}
\end{itemize}

\section{Determinantes}
É um valor numérico gerado por uma matriz. Ele revela, de forma bem resumida, propriedades geométricas e algébricas da matriz
\\[5pt]
\textbf{Para uma matriz possuir determinante, ela deve necessariamente ser uma matriz quadrada}
\\[5pt]
A notação que indica que estamos calculando o determinantes de uma matriz são duas barras verticais. Assim:
\[ A=
\left|
\begin{bmatrix}
    a & b \\
    c & d 
\end{bmatrix}
\right|
\] isso indica que estamos calculando o determinantes da matriz $A$

\subsection{Calculando Determinantes}
\begin{itemize}
    \item \textbf{Para matrizes 1x1 (unitária):} O determinante é o único valor da matriz (óbvio)
    \item \textbf{Para matrizes 2x2:} O determinante é obtido multiplicando as diagonais da matriz, e subtraindo o valor obtido multiplicando a diagonal principal ($\searrow$), pelo valor obtido multiplicando a diagonal secundária
    \item \textbf{Para matrizes 3x3:} Utilizamos a \textbf{Regra de Sarrus}
    \begin{itemize}
        \item \textbf{Primeiro passo:} Copiar as 2 primeiras colunas da matriz no fim dela
        \item \textbf{Segundo passo:} Somar os produtos das diagonais principais ($\searrow$) e subtrair os produtos das diagonais secundárias ($\swarrow$)
        \item \textbf{Fórmula:} Seja a matriz:
        \[
        A = \begin{bmatrix}
            a & b & c \\
            d & e & f \\
            g & h & i \\
        \end{bmatrix},
        \]
        o determinante é:
        \[
        \det(A) = aei + bfg + cdh - ceg - bdi - afh
        \]
    \end{itemize}
    
    \item \textbf{Para matrizes de ordem superior (4x4 em diante):} Utilizamos o \textbf{Método de Laplace}
    \begin{itemize}
        \item \textbf{O menor complementar:} O menor complementar relativo a um elemento é o determinante calculado ignorando a linha e a coluna do elemento relativo ao qual estamos calculando
        \item \textbf{O cofator:} O cofator \( C_{ij} \) de um elemento de uma matriz é dado por:
        \[
        C_{ij} = (-1)^{i+j} \cdot D_{ij}
        \]
        onde:
        \begin{itemize}
            \item \( D_{ij} \) é o menor complementar, ou seja, é o determinante da matriz obtida removendo a linha \( i \) e a coluna \( j \),
            \item \( (-1)^{i+j} \) determina o sinal do cofator
        \end{itemize} 
    \end{itemize}
\end{itemize}

\subsection{Propriedades dos Determinantes}

\section{Polinômios}
\subsection{Forma Geral de um Polinômio}
A raiz de um polinômio é número que faz ele dar zero

\section{Módulos}
Módulos são valores absolutos, distâncias
\subsection{Função Modular}
Módulo de X: Espelhamento horizontal
Módulo de Y (ou f(x)): Função nunca desce pra baixo do eixo y
\\[10pt]
Propriedade importante: $\sqrt{a^2} = |a|$
\subsection{Equações Modulares}
Esse tópico inteiro se baseia nas seguintes propriedades:
\[
|A| = B \Rightarrow A=B \text{ ou } A = -B
\]
Isso quer dizer que: Por ter um módulo, $x$ pode assumir 2 valores que satisfazem a igualdade
\\[5pt]
\textbf{Exemplos:}
\\[5pt]
$|3x - 2| = 7$
\\[5pt]
Logo
\\[5pt]
$3x - 2 = 7 \text{ ou } 3x - 2 = -7$
\\[5pt]
Resolvendo
\\[5pt]
$3x - 2 = 7 \Rightarrow 3x = 9 \Rightarrow x = 3$
\\
$3x - 2 = -7 \Rightarrow 3x = -5 \Rightarrow x = -\frac{5}{3}$
\subsection{Inequações Modulares}
NADA A VER COM INEQUAÇÕES, COLOQUEI ESSA IMAGEM AQUI POR PREGUIÇA
\\[10pt]
Uma forma de entender frações meio diferente é essa
\begin{figure}[h]
    \centering
    \includegraphics[width=0.5\linewidth]{gagagagagagagagagagagag.png}
    \caption{Número de fatias // Tamanho de 1 única fatia}
    \label{fig:enter-label}
\end{figure}
\end{document}
